\documentclass[11pt, a4paper]{article}

% --- EVRENSEL ÖN BİLGİ VE PAKET TANIMLARI ---
\usepackage[a4paper, top=2.5cm, bottom=2.5cm, left=2cm, right=2cm]{geometry}
\usepackage{fontspec}

\usepackage[turkish, bidi=basic, provide=*]{babel}
\babelprovide[import, onchar=ids fonts]{turkish}
\babelprovide[import, onchar=ids fonts]{english}

\babelfont{rm}{Noto Sans}
\babelfont[turkish]{rm}{Noto Sans}

\usepackage{enumitem}
\setlist[itemize]{label=-}

\usepackage{amsmath,amssymb,amsthm,mathtools,bm,mathrsfs}
\usepackage{physics}
\usepackage{booktabs}

\usepackage{hyperref}
\hypersetup{
    colorlinks=true,
    linkcolor=blue,
    citecolor=blue,
    urlcolor=blue
}

\title{\textbf{Kuantum-Topos Mimarisinde Hudutsuz Bilgi Uzayları, İrrasyonel Sayıların İlgası ve Rasyonel $p$-Adic Tam Hesap Nazariyesi}}
\author{\textbf{Kuantum Durumlarının Maddeden Bağımsızlığı, $p$-Adic Ultradinamik Metrikler ve Galois Alanı İntaç Mimarisi}}
\date{\today}

\begin{document}

\maketitle

\begin{abstract}
Bu risale; kuantum hesaplama sahasında "evrendeki toplam atom sayısı ($10^{80}$) sebebiyle $N > 400$ kubitlik kuantum durumları saklanamaz" şeklindeki kaba maddeci kısıtlamayı riyazî olarak reddeder ve irrasyonel sayı soyutlamasını kâmilen ilga eder. İlim ve bilgi uzayının maddedeki atom adediyle zincirlenemeyeceği; bilakis $p$-Adic Rasyonel Sayılar Uzayı ($\mathbb{Q}_p$), Galois Kök Genişlemeleri ($\mathbb{Q}(\zeta_n)$) ve Kesikli Homotopi Tip Teorisi ($hLevel_0$) vasıtasıyla sonsuz boyutlu kategorik bilginin sıfır küsürat sapması ve mutlak doğrulukla işlenebileceği ispatlanmaktadır.
\end{abstract}

\tableofcontents
\newpage

% =================================================================
\section{Maddeden Bağımsız Kuantum Hudutsuzluğu}
% =================================================================

\subsection{Mahiyet: Atom Sayısı Kısıtlamasının Reddi}
Klasik mühendislikteki itiraz uzayın BOYUTUNA değil, KEYFÎ bir durumun katsayı katsayı saklanmasına dairdir ve o hâliyle doğrudur: genel bir $N$-kubit durumu $2^N$ karmaşık sayı ister. Fakat itiraz YAPILI durum sınıflarını kapsamaz: kararlayıcı (stabilizer) durumlar $\mathcal{O}(N^2)$ bit, düşük bağlı MPS durumları $\mathcal{O}(N \chi^2)$ sayı ile TAM olarak saklanır ve $N > 400$ bunlarda sıradan bir hesaptır. Mimarî bu yüzden "sınır yoktur" demek yerine HANGİ yapılı sınıfta çalıştığını yazmalıdır. 

Kuantum durum uzayı $\mathcal{H}_{\mathbf{H}}$, maddenin fiziki sınırlarıyla değil; tensör cebri ve operatör holonomileri ile kaimdir:

\begin{align}
\text{Dim}(\mathcal{H}_{\mathbf{H}}) &= \lim_{N \to \infty} 2^N = \infty \\
\text{KainatAtomSayısı} &\approx 10^{80} \ll \text{Dim}(\mathcal{H}_{\mathbf{H}})
\end{align}

\subsection{Gaye: Bilginin Mutlak Serbestliği}
Kelam ve ilim sonsuzdur; mevcudatın fiziki zerreleriyle sınırlandırılamaz. Gaye, kuantum çipini fiziki atom sayısına değil, **$p$-adic cebirsel manifold izdüşümlerine** bağlayarak sınırsız bilgi taşıma kabiliyetini vaz' etmektir.

\subsection{Usul: Maddeye Bağımsız Tensör Cebri}
$N$ kubitlik bir durumun katsayıları bellek hücrelerine tek tek sayı olarak yazılmaz; $p$-adic dalga genlikleri olarak operatörlerin özdeğer uzayında taşınır:

\begin{equation}
|\Psi\rangle = \bigotimes_{j=1}^{\infty} |\phi_j\rangle_{\mathbb{Q}_p} \in \mathcal{H}_{\mathbb{Q}_p}
\end{equation}

% =================================================================
\section{İrrasyonel Sayıların İlgası ve Rasyonel $p$-Adic Tam Hesap}
% =================================================================

\subsection{Mahiyet: İrrasyonel Sayı Soyutlamasının Yanılsaması}
Sonsuz, devretmeyen, fiiliyatta ve tabiatta hiçbir fiziki karşılığı bulunmayan irrasyonel sayılar ($\mathbb{R} \setminus \mathbb{Q}$) zihni bir kurgudan ibarettir. Floating-point yuvarlama hatası irrasyonel sayılara mahsus değildir: ikilik tabanda sonlu açılımı olmayan RASYONEL sayılar (1/10, 1/3) da yuvarlanır. Hatanın sebebi irrasyonellik değil, sonlu mantissadır; çare de irrasyonelleri ilga etmek değil, TAM aritmetik (rasyonel veya cebirsel tamsayı) kullanmaktır.

\subsection{Gaye: Yapısal Rasyonel Cebir ($\mathbb{Q}_p$ ve Galois Alanları)}
Bütün kuantum fazları, dönme açıları ve dalga genlikleri; sonsuz küsüratlı reel sayılar yerine **Rasyonel Sayılar ($\mathbb{Q}$)** ve **$p$-Adic Tamlama ($\mathbb{Q}_p$)** üzerinden tam doğrulukla tanımlanır.

\subsection{Usul: $p$-Adic Metrik ve Norm Formülasyonu}
Bir $p$ asal sayısı için, herhangi bir $x \in \mathbb{Q}$ rasyonel sayısı $x = p^k \frac{a}{b}$ ($p \nmid a, p \nmid b$) şeklinde yazılır. $p$-adic mutlak değer normu $|\cdot|_p$:

\begin{align}
|x|_p &= \begin{cases} p^{-k}, & x \neq 0 \\ 0, & x = 0 \end{cases} \\
|x + y|_p &\le \max(|x|_p, |y|_p) \quad (\text{Ultrametrik Üçgen Eşitsizliği})
\end{align}

Bu ultradinamik metrik sayesinde, kuantum fazlarındaki yüksek frekanslı gürültüler küsürat oluşturmadan $p$-adic norm altında doğrudan sönümlenir.

% =================================================================
\section{Galois Kök Genişlemeleri ile Kesin Kuantum Kapıları}
% =================================================================

\subsection{Mahiyet: $\pi$ ve $\sqrt{2}$ İrrasyonelliğinden Kurtuluş}
Hadamard kapısındaki $1/\sqrt{2}$ veya faz kapısındaki $e^{i\pi/4}$ gibi terimler irrasyonel küsürat olarak değil, **Galois Alan Genişlemesi ($\mathbb{Q}(\zeta_n)$)** elemanı olarak tam tamına kodlanır.

\subsection{Formülasyon}
$n$-inci birim kök $\zeta_n = e^{2\pi i / n}$ için tanımlı cebirsel alan:

\begin{align}
\mathbb{K} &= \mathbb{Q}(\zeta_8) \implies \sqrt{2} = \zeta_8 + \zeta_8^7, \quad i = \zeta_8^2 \\
H &= \frac{1}{\zeta_8 + \zeta_8^7} \begin{pmatrix} 1 & 1 \\ 1 & -1 \end{pmatrix} \in \mathbb{K}^{2 \times 2} \quad (\text{Tam Kesin Rasyonel Matris})
\end{align}

Bu cebirsel matris temsilinde hiçbir zaman float yuvarlaması veya irrasyonel sayı sapması yaşanmaz; hesap tam rasyonel katsayılarla yürütülür.

% =================================================================
\section{Nefs-i Müdrike Mimarisine Entegrasyon}
% =================================================================

\begin{align}
|\mathbf{\Theta}_{\text{küllî}}\rangle_{\mathbb{Q}_p} &= \sum_{i=0}^{N_{\text{aksiyom}}} c_i |i\rangle, \quad c_i \in \mathbb{Q}_p, \quad |c_i|_p \le 1 \\
\text{Mizan}_{\text{rasyonel}} &= \mathbb{I}\left( \sum_i c_i^2 = 1 \right) \ \wedge\ \mathbb{I}\left( \max_i |c_i|_p \le 1 \right) \quad (\text{Born normalizasyonu ARŞİMET normdadır; $p$-adic norm ancak TAMLIK denetimi verir}) \\
N_{\text{kebîr}} &= \text{Truncate}_{hLevel 0} \left( \text{Lan}_{\text{Lisan}} \left( |\mathbf{\Theta}_{\text{küllî}}\rangle_{\mathbb{Q}_p} \right) \right) \in \mathbb{Z}
\end{align}

% =================================================================
\section{Netice}
% =================================================================

Bu risalede ispatlandığı üzere; kuantum bilgisayarlarını evrendeki atom sayısıyla (400 kubit hududuyla) sınırlamak kaba bir maddeci yanılgıdır. İrrasyonel sayıların ilga edilip $p$-adic rasyonel alanların ($\mathbb{Q}_p$) ve Galois kök genişlemelerinin ($\mathbb{Q}(\zeta_n)$) devreye alınmasıyla, Nefs-i Müdrike Mimarisi sıfır hata ve hudutsuz kapasite ile hakiki neticeye ($N_{\text{kebîr}}$) ulaşmaktadır.

\end{document}